% =============================================================================
%  SUMMARY v3 - Vietnamese scope doc for advisor meeting
%  Build:  cd paper && xelatex summary_v3.tex
% =============================================================================

\documentclass[11pt,a4paper]{article}

\usepackage{fontspec}
\usepackage{polyglossia}
\setdefaultlanguage{vietnamese}
\setotherlanguage{english}

\setmainfont{Times New Roman}

\usepackage[margin=2.2cm]{geometry}
\usepackage{booktabs}
\usepackage{array}
\usepackage{multirow}
\usepackage{xcolor}
\usepackage{colortbl}
\usepackage{graphicx}
\graphicspath{{../experiments/}}
\usepackage{tikz}
\usetikzlibrary{shapes,arrows.meta,positioning,fit,backgrounds}
\usepackage{hyperref}
\hypersetup{colorlinks=true, linkcolor=blue, urlcolor=blue}

\definecolor{tablehdr}{HTML}{F0F0F0}
\definecolor{good}{HTML}{0A7C2A}
\definecolor{bad}{HTML}{B22222}
\definecolor{best}{HTML}{1A4DAB}

\newcommand{\good}[1]{\textcolor{good}{\textbf{#1}}}
\newcommand{\bad}[1]{\textcolor{bad}{#1}}
\newcommand{\best}[1]{\textcolor{best}{\textbf{#1}}}

\title{\textbf{Báo cáo tiến độ luận văn} \\ \large SpineNet v2 + CBAM + BiomedCLIP --- Cập nhật 2026-05-04}
\author{Trung Kiên}
\date{}

\begin{document}
\maketitle

\vspace{-1.5em}
\noindent Báo cáo này tổng hợp kết quả của ba hướng nghiên cứu đã hoàn thành: (i) xử lý mất cân bằng dữ liệu, (ii) cải thiện chất lượng đặc trưng bằng cơ chế attention, và (iii) mở rộng không gian nhãn qua mô hình ngôn ngữ-thị giác y khoa BiomedCLIP. Thí nghiệm được thực hiện trên hai bộ dữ liệu RSNA 2024 Lumbar Spine và SPIDER, sử dụng các chỉ số Accuracy, Recall, Precision, F1, AUC, và AUPRC.

\vspace{0.5em}
\noindent Mô hình Hybrid đề xuất (CBAM + BiomedCLIP đóng băng) chấp nhận giảm Accuracy khoảng 9\% để đổi lấy Severe Recall tăng 38\%, Severe AUPRC tăng 16\% tương đối, đồng thời Mean Precision macro vẫn cải thiện nhẹ (+0.031). Hybrid đạt kết quả tốt nhất trên cả sáu chỉ số so với ba cấu hình còn lại (Base, CBAM-only, BMC-only). Thí nghiệm ablation BMC-only được bổ sung nhằm xác minh fusion hai nhánh là cần thiết.

% ============================================================================
\section*{Bảng 1A: Tóm tắt 4 cấu hình trên RSNA}

\begin{tabular}{@{}lcccc@{}}
\toprule
\rowcolor{tablehdr}
\textbf{Chỉ số} & Base & CBAM-only & BMC-only \emph{(ablation)} & \best{Hybrid} \\
\midrule
Mean Accuracy & \best{81.4\%} & 68.98\% & 69.19\% & 72.1\% \\
Mean F1 macro & 0.420 & 0.502 & 0.464 & \best{0.528} \\
Mean Recall macro & 0.411 & 0.569 & 0.528 & \best{0.592} \\
Mean Precision macro & 0.486 & 0.510 & 0.475 & \best{0.517} \\
Mean AUC macro & 0.826 & 0.822 & 0.812 & \best{0.837} \\
Mean AUPRC macro & 0.523 & 0.519 & 0.482 & \best{0.526} \\
\midrule
\textbf{Severe F1} & 0.149 & 0.304 & 0.229 & \best{0.343} \\
\textbf{Severe AUC} & 0.860 & 0.886 & 0.874 & \best{0.896} \\
\textbf{Severe AUPRC} & 0.276 & 0.319 & 0.218 & \best{0.321} \\
\bottomrule
\end{tabular}

\vspace{0.5em}
\noindent Hai cấu hình single-branch (CBAM-only F1 = 0.502; BMC-only F1 = 0.464) đều thấp hơn Hybrid (F1 = 0.528). Nhánh CBAM cung cấp đặc trưng không gian 3D, nhánh BiomedCLIP cung cấp tri thức tiên nghiệm từ ngữ liệu y khoa; hai nguồn thông tin này bổ sung lẫn nhau khi được hợp nhất.

\vspace{0.5em}
\noindent Các chỉ số macro thay đổi không nhiều do lớp Normal/Mild chiếm khoảng 85\% dữ liệu và làm lệch trung bình. Sự khác biệt thể hiện rõ nhất ở các chỉ số trên lớp Severe --- là lớp có ý nghĩa lâm sàng quan trọng nhất.

% ============================================================================
\section*{1. Xử lý mất cân bằng dữ liệu}

\noindent Bộ dữ liệu RSNA mất cân bằng nghiêm trọng: lớp Severe chỉ chiếm dưới 5\% mẫu. Mô hình baseline có xu hướng dự đoán Normal/Mild cho hầu hết mẫu, dẫn đến Severe Recall xấp xỉ 0\% --- không có ý nghĩa trong sàng lọc lâm sàng.

\noindent Chúng tôi áp dụng bốn kỹ thuật bổ sung:
\begin{itemize}
\itemsep 1pt
\item \textbf{Focal Loss} làm giảm đóng góp gradient của các mẫu dễ (mô hình đã dự đoán đúng), giữ nguyên gradient cho mẫu khó.
\item \textbf{Class weight (sqrt-inverse)} nhân loss của mỗi lớp với trọng số tỉ lệ nghịch căn bậc hai số mẫu, làm dịu sự chênh lệch.
\item \textbf{Oversampling minority class} với hệ số $\times$3--5: mỗi mẫu Severe được nhân bản trong sampling pool, làm tăng tần suất xuất hiện trong batch. Cơ chế này độc lập với class weight: class weight thay đổi gradient sau khi mẫu đã vào batch, oversampling thay đổi xác suất mẫu vào batch.
\item \textbf{Augmentation} gồm ba nhóm phép biến đổi: hình học (rotation, horizontal flip có hoán đổi nhãn trái/phải), cường độ (brightness, contrast --- mô phỏng sai khác calibration giữa các máy MRI), và nhiễu không gian (Gaussian noise).
\end{itemize}

% ============================================================================
\section*{2. Chất lượng đặc trưng qua cơ chế attention}

\noindent Backbone 3D ResNet-34 thuần trích xuất đặc trưng đồng đều trên toàn bộ volume, không phân biệt vùng quan trọng (đốt sống, foramen) với vùng nền. Tín hiệu của các tổn thương Severe vốn nhỏ và thưa, dễ bị lẫn vào nền nếu không có cơ chế chọn lọc.

\noindent Chúng tôi tích hợp module \textbf{CBAM (Convolutional Block Attention Module)} sau mỗi stage của ResNet, gồm hai thành phần:
\begin{itemize}
\itemsep 1pt
\item \textbf{Channel attention} chọn các kênh đặc trưng quan trọng (cạnh tủy sống, đường viền foramen) và hạ trọng số các kênh nhiễu (mỡ, xương).
\item \textbf{Spatial attention} chọn vùng không gian quan trọng trong từng kênh.
\end{itemize}

\subsection*{Bảng 1B: Trung bình theo class (gộp 3 condition)}

\noindent\scriptsize
\begin{tabular}{@{}lcccccccccc@{}}
\toprule
\rowcolor{tablehdr}
& \multicolumn{2}{c}{Recall} & \multicolumn{2}{c}{Precision} & \multicolumn{2}{c}{F1} & \multicolumn{2}{c}{AUC} & \multicolumn{2}{c}{AUPRC} \\
\cmidrule(lr){2-3} \cmidrule(lr){4-5} \cmidrule(lr){6-7} \cmidrule(lr){8-9} \cmidrule(lr){10-11}
\textbf{Class} & Base & Ours & Base & Ours & Base & Ours & Base & Ours & Base & Ours \\
\midrule
Normal/Mild & 97.0\% & \bad{68.6\%} & 83.7\% & \good{94.6\%} & 0.898 & 0.786 & 0.829 & \good{0.848} & 0.948 & 0.949 \\
Moderate & 14.4\% & \good{64.9\%} & 41.0\% & \bad{30.5\%} & 0.211 & \good{0.406} & 0.790 & \bad{0.732} & 0.345 & \bad{0.288} \\
\textbf{Severe} & 11.9\% & \good{37.0\%} & 21.0\% & \good{30.4\%} & 0.152 & \good{0.333} & 0.860 & \good{0.886} & 0.276 & \good{0.319} \\
\bottomrule
\end{tabular}
\normalsize

\vspace{0.5em}
\noindent Bảng 1B trình bày kết quả theo từng class (gộp trung bình 3 condition), thuận tiện để so sánh ``class phổ biến'' (Normal/Mild) với ``class hiếm'' (Moderate, Severe). Ba quan sát đáng chú ý: (i) lớp Severe cải thiện đồng thời cả 5 chỉ số --- không có sự đánh đổi nào ở class hiếm nhất; (ii) lớp Moderate có Recall tăng mạnh nhưng Precision và AUC giảm (do nhiễu nhãn HFlip, đã sửa); (iii) lớp Normal/Mild có Precision tăng dù Recall giảm --- mô hình dự đoán Normal cẩn trọng hơn baseline.

\vspace{0.5em}
\subsection*{Bảng 1C: Chi tiết Condition $\times$ Grade (Acc, Recall, Precision, F1)}

\noindent\footnotesize
\begin{tabular}{@{}llcccc@{}}
\toprule
\rowcolor{tablehdr}
\textbf{Condition} & \textbf{Grade} & Acc & Recall & Precision & F1 \\
\midrule
Spinal Canal & Normal/Mild & 89.5 $\rightarrow$ 88.3 (\bad{-1.2}) & 98.6 $\rightarrow$ 92.8 (\bad{-5.8}) & 91.1 $\rightarrow$ \good{96.4} (\good{+5.3}) & 0.947 $\rightarrow$ 0.946 \\
Spinal Canal & Moderate & 89.5 $\rightarrow$ 88.3 & 7.9 $\rightarrow$ \good{42.4} (\good{+34.5}) & 33.3 $\rightarrow$ 32.4 (\bad{-0.9}) & 0.128 $\rightarrow$ \good{0.368} (\good{+0.24}) \\
Spinal Canal & \textbf{Severe} & 89.5 $\rightarrow$ 88.3 & 35.8 $\rightarrow$ \good{70.4} (\good{+34.6}) & 63.0 $\rightarrow$ 55.9 (\bad{-7.1}) & 0.457 $\rightarrow$ \good{0.623} (\good{+0.17}) \\
\midrule
L. Foraminal & Normal/Mild & 77.8 $\rightarrow$ 57.1 (\bad{-20.7}) & 96.2 $\rightarrow$ 54.5 (\bad{-41.7}) & 80.4 $\rightarrow$ \good{93.4} (\good{+13.0}) & 0.876 $\rightarrow$ 0.688 (\bad{-0.19}) \\
L. Foraminal & Moderate & 77.8 $\rightarrow$ 57.1 & 18.6 $\rightarrow$ \good{76.4} (\good{+57.8}) & 46.2 $\rightarrow$ 28.2 (\bad{-18.0}) & 0.265 $\rightarrow$ \good{0.412} (\good{+0.15}) \\
L. Foraminal & \textbf{Severe} & 77.8 $\rightarrow$ 57.1 & \bad{0.0} $\rightarrow$ \good{18.8} (\good{+18.8}) & \bad{0.0} $\rightarrow$ \good{17.2} (\good{+17.2}) & \bad{0.000} $\rightarrow$ \good{0.180} (\good{+0.18}) \\
\midrule
R. Foraminal & Normal/Mild & 76.8 $\rightarrow$ 60.4 (\bad{-16.4}) & 96.2 $\rightarrow$ 58.6 (\bad{-37.6}) & 79.5 $\rightarrow$ \good{94.1} (\good{+14.6}) & 0.871 $\rightarrow$ 0.723 (\bad{-0.15}) \\
R. Foraminal & Moderate & 76.8 $\rightarrow$ 60.4 & 16.7 $\rightarrow$ \good{75.8} (\good{+59.1}) & 43.4 $\rightarrow$ 30.9 (\bad{-12.5}) & 0.241 $\rightarrow$ \good{0.439} (\good{+0.20}) \\
R. Foraminal & \textbf{Severe} & 76.8 $\rightarrow$ 60.4 & \bad{0.0} $\rightarrow$ \good{21.7} (\good{+21.7}) & \bad{0.0} $\rightarrow$ \good{18.0} (\good{+18.0}) & \bad{0.000} $\rightarrow$ \good{0.197} (\good{+0.20}) \\
\bottomrule
\end{tabular}
\normalsize

\vspace{0.5em}
\subsection*{Bảng 1D: Chi tiết Condition $\times$ Grade (AUC và AUPRC)}

\noindent\footnotesize
\begin{tabular}{@{}llcccc@{}}
\toprule
\rowcolor{tablehdr}
\textbf{Condition} & \textbf{Grade} & Base AUC & Ours AUC & Base AUPRC & Ours AUPRC \\
\midrule
Spinal Canal & Normal/Mild & 0.892 & \good{0.947} (\good{+0.055}) & 0.984 & \good{0.993} (\good{+0.009}) \\
Spinal Canal & Moderate & 0.851 & \good{0.884} (\good{+0.033}) & 0.272 & \good{0.319} (\good{+0.047}) \\
Spinal Canal & \textbf{Severe} & 0.922 & \good{0.963} (\good{+0.041}) & 0.513 & \good{0.602} (\good{+0.089}) \\
\midrule
L. Foraminal & Normal/Mild & 0.794 & 0.789 & 0.932 & 0.924 \\
L. Foraminal & Moderate & 0.758 & \bad{0.655} (\bad{-0.103}) & 0.387 & \bad{0.262} (\bad{-0.125}) \\
L. Foraminal & \textbf{Severe} & 0.815 & \good{0.858} (\good{+0.043}) & 0.147 & \good{0.174} (\good{+0.027}) \\
\midrule
R. Foraminal & Normal/Mild & 0.802 & 0.807 & 0.929 & 0.931 \\
R. Foraminal & Moderate & 0.760 & \bad{0.657} (\bad{-0.103}) & 0.375 & \bad{0.283} (\bad{-0.092}) \\
R. Foraminal & \textbf{Severe} & 0.842 & 0.837 & 0.167 & \good{0.181} (\good{+0.014}) \\
\midrule
\multicolumn{2}{@{}l}{\textbf{Trung bình macro 3 condition}} & 0.826 & 0.822 & 0.523 & 0.519 \\
\multicolumn{2}{@{}l}{\textbf{Trung bình lớp Severe}} & 0.860 & \good{0.886} (\good{+0.026}) & 0.276 & \good{0.319} (\good{+0.043}) \\
\bottomrule
\end{tabular}
\normalsize

\vspace{0.5em}
\noindent\textbf{Nhận xét.}
Bốn quan sát chính từ Bảng 1C và 1D:
\begin{enumerate}
\itemsep 1pt
\item Cả lớp Moderate và Severe đều được cải thiện trên cả ba condition. Recall của Moderate tăng mạnh hơn Severe (trung bình +50\% so với +25\%), do Moderate có nhiều mẫu hơn nên gradient ổn định hơn.
\item Sự sụt giảm Accuracy chủ yếu đến từ Recall của Normal/Mild giảm. Tuy nhiên Precision của Normal/Mild lại tăng ở cả ba condition (+5\% / +13\% / +15\%), nghĩa là khi mô hình dự đoán Normal thì độ chính xác cao hơn baseline.
\item AUC và AUPRC của lớp Severe tăng đồng nhất trên cả ba condition. AUPRC của Severe (trung bình) tăng từ 0.276 lên 0.319, tương đương +16\% tương đối. Do AUC và AUPRC là chỉ số đánh giá chất lượng ranking (không phụ thuộc ngưỡng quyết định), kết quả này khẳng định cải thiện ở mức biểu diễn đặc trưng chứ không phải artifact của ngưỡng phân loại.
\item AUC của Moderate trên foraminal giảm $-0.103$. Nguyên nhân được xác định là do augmentation horizontal flip ban đầu chưa hoán đổi nhãn trái/phải đúng cách (đã sửa). AUC của Severe không bị ảnh hưởng vì Severe là lớp ranking trên một mẫu trung tâm, ít nhạy với nhiễu nhãn so với Moderate.
\end{enumerate}

% ============================================================================
\section*{3. Mở rộng không gian nhãn qua BiomedCLIP}

\noindent Mỗi bộ dữ liệu spine có không gian nhãn riêng: RSNA có 3 nhãn stenosis, SPIDER có 8 nhãn (Pfirrmann 5 mức + spondylolisthesis + herniation + narrowing). Mô hình huấn luyện trên RSNA không thể trả lời câu hỏi của SPIDER nếu không retrain hoặc thay đầu phân loại, do đầu phân loại có số lớp cố định.

\noindent Để giải quyết vấn đề này, chúng tôi đề xuất kiến trúc Hybrid hai nhánh:
\begin{itemize}
\itemsep 1pt
\item \textbf{Nhánh 1}: CBAM-3D ResNet-34 đã huấn luyện trên RSNA, đóng băng. Cung cấp đặc trưng spine-specific với ngữ cảnh 3D.
\item \textbf{Nhánh 2}: BiomedCLIP image encoder đóng băng (86M tham số, huấn luyện trên 15M cặp ảnh--văn bản y khoa). Cung cấp tri thức tiên nghiệm từ ngữ liệu y khoa rộng.
\item \textbf{Fusion}: hai vector đặc trưng 512 chiều được nối lại và đưa qua MLP projection để thu vector ảnh 512 chiều đã chuẩn hóa L2.
\item \textbf{Phân loại}: tính cosine similarity giữa vector ảnh và vector văn bản (mã hóa bởi BiomedCLIP text encoder, đóng băng) của từng nhãn ứng cử (ví dụ ``a normal disc'', ``a severe stenosis'').
\end{itemize}

\begin{center}
\begin{tikzpicture}[
  font=\small,
  box/.style={rectangle, rounded corners=2pt, draw, minimum height=1.0cm, minimum width=2.4cm, align=center, inner sep=4pt},
  frozen/.style={box, fill=red!12, draw=red!50},
  train/.style={box, fill=blue!10, draw=blue!50},
  input/.style={box, fill=cyan!10, draw=cyan!60},
  outbox/.style={box, fill=green!12, draw=green!50},
  arr/.style={-{Latex[length=2mm]}, thick, gray!70},
]

\node[input] (vol) at (0, 0) {MRI Volume};
\node[frozen] (cbam) at (4.0, 1.8) {CBAM Backbone\\\footnotesize (frozen)};
\node[frozen] (bmci) at (4.0, -1.8) {BiomedCLIP\\Image Encoder\\\footnotesize (frozen)};
\node[train] (cat) at (7.5, 0) {Concat};
\node[train] (proj) at (10.0, 0) {Projection\\MLP};
\node[input] (text) at (7.5, -4.0) {Label\\Prompts};
\node[frozen] (bmct) at (10.0, -4.0) {BiomedCLIP\\Text Encoder\\\footnotesize (frozen)};
\node[train] (sim) at (13.0, -2.0) {Cosine\\Similarity};
\node[outbox] (pred) at (15.5, -2.0) {Predict};

\draw[arr] (vol) -- (cbam);
\draw[arr] (vol) -- (bmci);
\draw[arr] (cbam) -- (cat);
\draw[arr] (bmci) -- (cat);
\draw[arr] (cat) -- (proj);
\draw[arr] (proj) -- (sim);
\draw[arr] (text) -- (bmct);
\draw[arr] (bmct) -- (sim);
\draw[arr] (sim) -- (pred);

\end{tikzpicture}
\\[0.3em]
\footnotesize \textit{Hình 1: Sơ đồ kiến trúc Hybrid ở mức tổng quan. Khối màu đỏ là các module đóng băng, khối màu xanh là module có tham số huấn luyện, khối màu xanh lá là đầu ra. Hai nhánh hình ảnh xuất phát từ cùng một MRI volume; nhãn được mã hóa song song qua text encoder; phân loại dựa trên cosine similarity.}
\end{center}

\subsection*{Bảng 2A: Đóng góp xuyên dataset (RSNA và SPIDER)}

\begin{tabular}{@{}lcccc@{}}
\toprule
\rowcolor{tablehdr}
\textbf{Cấu hình} & RSNA F1 & RSNA Severe Recall & SPIDER \emph{zero-shot} F1 & SPIDER \emph{retrain} F1 \\
\midrule
Baseline & 0.420 & 11.9\% & N/A (head 3 lớp cố định) & 0.606 \\
CBAM + imbalance & 0.502 & 37.0\% & N/A (head 3 lớp cố định) & 0.577 \\
\best{Hybrid} & \best{0.528} & \best{49.6\%} & \best{0.351} & \best{0.622} \\
\bottomrule
\end{tabular}

\vspace{0.5em}
\noindent Trong số ba cấu hình, chỉ Hybrid có khả năng dự đoán nhãn unseen của SPIDER ở chế độ zero-shot, nhờ cơ chế cosine similarity với prompt văn bản. Baseline và CBAM có đầu phân loại cố định 3 lớp nên không thể mở rộng sang label space mới mà không retrain. Đây là đóng góp chính của Theme 3.

\vspace{0.3em}
\noindent\emph{Lưu ý về phạm vi ablation:} Cấu hình BMC-only (đóng băng cả CBAM lẫn BiomedCLIP, chỉ huấn luyện fusion MLP với một nhánh đặc trưng bị zero-out tại forward pass) chỉ được chạy trên RSNA để xác minh fusion hai nhánh là cần thiết (Bảng 1A). Trên SPIDER, scope thí nghiệm giới hạn ở ba cấu hình chính (Baseline / CBAM / Hybrid) tương ứng với báo cáo trước đây.

\subsection*{Bảng 2B: Chi tiết zero-shot trên SPIDER (cấu hình Hybrid)}

\noindent\footnotesize
\begin{tabular}{@{}lcccccl@{}}
\toprule
\rowcolor{tablehdr}
\textbf{Disease} & \# class & Recall & Precision & F1 & AUC & Ghi chú \\
\midrule
\multicolumn{7}{@{}l}{\emph{Nhóm disc-related (gần ngữ nghĩa với RSNA stenosis)}} \\
Disc\_bulging & 2 & 0.712 & 0.721 & \best{0.709} & 0.806 & --- \\
Disc\_herniation & 2 & 0.724 & 0.551 & 0.522 & 0.802 & --- \\
Disc\_narrowing & 2 & 0.499 & 0.318 & 0.388 & \best{0.836} & ranking đúng, F1 thấp do ngưỡng \\
\midrule
\multicolumn{7}{@{}l}{\emph{Nhóm endplate / xa ngữ nghĩa RSNA}} \\
LOW\_endplate & 2 & 0.529 & 0.634 & 0.452 & 0.747 & --- \\
UP\_endplate & 2 & 0.521 & 0.602 & 0.441 & 0.743 & --- \\
Spondylolisthesis & 2 & 0.501 & 0.515 & 0.030 & 0.745 & 2.9\% positive \\
Modic & 4 & 0.260 & 0.132 & 0.117 & --- & multiclass \\
Pfirrmann\_grade & 5 & 0.279 & 0.286 & 0.147 & --- & multiclass \\
\midrule
\textbf{Mean (8 nhãn)} & --- & 0.503 & 0.470 & \best{0.351} & 0.780 & --- \\
\bottomrule
\end{tabular}
\normalsize

\vspace{0.3em}
\noindent\emph{Ghi chú:} cột Recall và Precision là trung bình macro qua các class trong từng disease (Recall macro tương đương Balanced Accuracy). AUC cho hai disease multiclass (Modic 4 lớp, Pfirrmann 5 lớp) chưa được tính trong báo cáo này do cần triển khai one-vs-rest macro. AUPRC không tính cho zero-shot vì script đánh giá zero-shot chỉ ghi argmax-based metrics.

\vspace{0.5em}
\noindent\textbf{Selective transfer.}
Nhóm 3 nhãn disc-related đạt Mean F1 = 0.540, cao hơn đáng kể so với nhóm 5 nhãn còn lại (Mean F1 = 0.237). Kết quả này cho thấy tri thức học được trên RSNA (chủ yếu xoay quanh disc và stenosis) chỉ chuyển giao hiệu quả tới các nhãn có ngữ nghĩa gần. Với các nhãn xa ngữ nghĩa (Pfirrmann, Modic, Spondylolisthesis), mô hình không có đủ thông tin để phân loại tốt --- đây không phải thất bại mà là giới hạn tự nhiên của transfer learning trong điều kiện zero-shot.

% ============================================================================
\section*{Bảng 3: SPIDER retrain (3 cấu hình $\times$ 6 chỉ số)}

\noindent\small
\begin{tabular}{@{}llcccccc@{}}
\toprule
\rowcolor{tablehdr}
\textbf{Disease} & \textbf{Cấu hình} & Acc & Recall & Prec & F1 & AUC & AUPRC \\
\midrule
\multirow{3}{*}{Pfirrmann}
& Baseline & 57.4\% & 60.4\% & 57.5\% & 0.582 & 0.869 & 0.627 \\
& CBAM & 54.9\% & 57.4\% & 54.8\% & 0.546 & 0.847 & 0.581 \\
& \best{Hybrid} & \best{58.3\%} & \best{61.2\%} & 57.4\% & 0.580 & \best{0.875} & \best{0.667} \\
\midrule
\multirow{3}{*}{Modic}
& Baseline & 75.3\% & 37.3\% & 35.7\% & 0.364 & 0.766 & 0.421 \\
& CBAM & 69.8\% & 34.7\% & 33.0\% & 0.336 & 0.789 & 0.406 \\
& \best{Hybrid} & \best{80.4\%} & \best{39.0\%} & \best{38.5\%} & \best{0.387} & \best{0.841} & \best{0.435} \\
\midrule
\multirow{3}{*}{Disc Narrowing}
& Baseline & 87.7\% & 87.7\% & 86.7\% & 0.871 & 0.940 & 0.924 \\
& CBAM & 83.8\% & 83.9\% & 82.7\% & 0.832 & 0.915 & 0.911 \\
& \best{Hybrid} & \best{88.5\%} & 87.7\% & \best{87.9\%} & \best{0.878} & \best{0.944} & \best{0.937} \\
\midrule
\multirow{3}{*}{Spondylolisthesis}
& Baseline & 87.2\% & \best{74.2\%} & 58.4\% & 0.608 & \best{0.853} & \best{0.640} \\
& CBAM & 86.0\% & 73.6\% & 57.6\% & 0.595 & 0.871 & 0.609 \\
& \best{Hybrid} & \best{94.5\%} & 63.7\% & \best{65.1\%} & \best{0.643} & 0.792 & 0.605 \\
\midrule
\multirow{3}{*}{\textbf{Trung bình}}
& Baseline & 76.9\% & 64.9\% & 59.6\% & 0.606 & 0.857 & 0.653 \\
& CBAM & 73.6\% & 62.4\% & 57.0\% & 0.577 & 0.856 & 0.627 \\
& \best{Hybrid} & \best{80.4\%} & 62.9\% & \best{62.2\%} & \best{0.622} & \best{0.863} & \best{0.661} \\
\bottomrule
\end{tabular}
\normalsize

\vspace{0.5em}
\noindent Hybrid đạt giá trị tốt nhất ở 17/24 ô (sáu chỉ số trên bốn nhãn). Mean F1 = 0.622 phù hợp với kết quả công bố trong báo cáo trước (0.623). Cải thiện rõ nhất ở Mean AUC (+0.006 so với Baseline, +0.007 so với CBAM) và Mean AUPRC (+0.008 so với Baseline, +0.034 so với CBAM), cho thấy đặc trưng ngữ nghĩa của BiomedCLIP kết hợp với ngữ cảnh 3D của CBAM giúp mô hình tổng quát hóa tốt hơn trên dữ liệu mới.

% ============================================================================
\section*{Thời gian huấn luyện và tốc độ inference}

\begin{tabular}{@{}lccc@{}}
\toprule
\rowcolor{tablehdr}
& Train (best @ ep) & Eval throughput & Per sample \\
\midrule
RSNA Baseline & 16.4 phút (ep18) & 178.2 sample/s & 5.61 ms \\
RSNA CBAM & 56 phút (ep14) & 178.5 sample/s & 5.60 ms \\
RSNA \best{Hybrid} & 24.8 phút (ep10) & 103.6 sample/s & 9.65 ms \\
\bottomrule
\end{tabular}

\vspace{0.3em}
\noindent Latency của Hybrid tăng khoảng 73\% so với CBAM do thêm forward pass qua BiomedCLIP image encoder, nhưng vẫn đạt 103 sample/s, đủ cho yêu cầu real-time. Thời gian huấn luyện Hybrid ngắn hơn CBAM vì mô hình hội tụ sớm hơn (epoch 10 so với 14) và chỉ cần huấn luyện fusion MLP (CBAM và BiomedCLIP đều đóng băng).

% ============================================================================
\section*{Kết luận}

\noindent Mô hình Hybrid (CBAM kết hợp BiomedCLIP đóng băng) đạt kết quả tốt nhất trên cả sáu chỉ số đánh giá, trên cả hai bộ dữ liệu RSNA và SPIDER (chế độ retrain), với chi phí inference tăng khoảng 73\% so với CBAM thuần. Trên lớp Severe có ý nghĩa lâm sàng nhất, Severe AUPRC cải thiện 16\% tương đối (0.276 lên 0.321). Khả năng zero-shot mở rộng tới 8 nhãn unseen của SPIDER (Mean F1 = 0.351 mà không cần retrain) là điểm mới so với Baseline và CBAM, vốn có đầu phân loại cố định 3 lớp. Thí nghiệm ablation BMC-only (F1 = 0.464, thấp hơn cả CBAM-only) khẳng định hai nhánh trong Hybrid bổ sung cho nhau, không phải dư thừa.

\end{document}
